%========================================================================================
% A) PACKAGES
%========================================================================================

\documentclass[journal,twoside,web]{ieeecolor}
\usepackage{generic}
\usepackage{cite}
\usepackage{amsmath,amssymb,amsfonts}
\usepackage{algorithmic}
\usepackage{graphicx}
\usepackage{textcomp}

 %Hyphenation
 \tolerance=1
 \emergencystretch=\maxdimen
 \hyphenpenalty=10000
 \hbadness=10000

%Figures
\usepackage{float}
\usepackage[titles,subfigure]{tocloft}
\usepackage{subfig}
\usepackage{fancyhdr}
\usepackage{hyperref}
\usepackage{rotating}
\usepackage[export]{adjustbox}

\usepackage{booktabs}%
\usepackage{MnSymbol,wasysym}
\usepackage{multirow}%
\usepackage{amsmath,amssymb,amsfonts}%
\usepackage{mathrsfs}%
\usepackage{textcomp}%
\usepackage{manyfoot}%
\usepackage{listings}%
\usepackage{MnSymbol,wasysym}

\definecolor{devyblue}{rgb}{0.2, 0.4, 0.6}



%==================================================
% B) Encabezado
%==================================================

    %--------------------------------------------------
    % B.1) Encabezado
    %--------------------------------------------------
    
\def\BibTeX{{\rm B\kern-.05em{\sc i\kern-.025em b}\kern-.08em
    T\kern-.1667em\lower.7ex\hbox{E}\kern-.125emX}}
\markboth{Ingeniería de Ejecución en Computación e Informática, Seminario de Computación e Informática, 1-2026}
{Vásquez-Salgado \MakeLowercase{\textit{et al.}}: Título de la propuesta}


    %--------------------------------------------------
    % B.2) Título de la propuesta
    %--------------------------------------------------

\begin{document}
\title{Análisis estructural y comparativo de programas de Ingeniería Informática en Chile basado en procesamiento de lenguaje natural y clustering}


    %--------------------------------------------------
    % B.3) Autores
    %--------------------------------------------------

    
\author{Fabián Ibarra Norambuena, Manuel Villalobos Cid
\thanks{Trabajo enviado el 23 de abril de 2026}
\thanks{Fabián Ibarra Norambuena es estudiante de Ingeniería de Ejecución en Computación e Informática (email: fabian.ibarra.n@usach.cl).}
\thanks{Manuel Villalobos Cid es académico del Departamento de Ingeniería Informática (email: manuel.villalobos@usach.cl).}
}

\maketitle

    %--------------------------------------------------
    % B.4) Resumen
    %--------------------------------------------------
    
\begin{abstract}
Este trabajo aborda la heterogeneidad existente en los programas de Ingeniería Informática en Chile mediante un enfoque de análisis computacional de mallas curriculares y descripciones de asignaturas. Se propone una metodología basada en procesamiento de lenguaje natural para la representación textual de los contenidos formativos, junto con técnicas de aprendizaje no supervisado, específicamente clustering, para identificar patrones estructurales entre programas. Adicionalmente, se incorporan métricas de similitud y análisis comparativo de perfiles de egreso y distribución de áreas disciplinares, con el fin de complementar el agrupamiento y proporcionar una caracterización más integral. El estudio considera datos recopilados desde fuentes institucionales públicas, permitiendo construir una visión sistemática y actualizada de la oferta académica. Como resultado, se espera evidenciar diferencias significativas en enfoques formativos, niveles de especialización y coherencia curricular entre instituciones, contribuyendo a una mejor comprensión del sistema y proporcionando una base objetiva para la comparación de programas.

\noindent \textbf{Categoría del trabajo:} Investigación (con matices de Desarrollo). \\
\textbf{Origen del tema:} Propuesto por profesor.
\end{abstract}



    %--------------------------------------------------
    % B.5) Keywords
    %--------------------------------------------------
\begin{IEEEkeywords}
Educación en Ingeniería, Algoritmos de agrupamiento, Procesamiento de lenguaje natural, Programas educativos, Minería de textos


\end{IEEEkeywords}

%==================================================
% C) Documento
%==================================================

    %--------------------------------------------------
    % C.1) Resumen de la propuesta
    %--------------------------------------------------

\section{Resumen de la propuesta}\label{sec:resumen}

        %--------------------------------------------------
        % C.1.1) Plantamiento del problema
        %--------------------------------------------------
    
\subsection{Planteamiento del problema} \label{planteamiento}

\IEEEPARstart{E}{l} vertiginoso avance tecnológico y la creciente digitalización han consolidado a la Ingeniería Informática como una disciplina fundamental en Chile, propiciando un aumento significativo en su demanda y en su impacto dentro del mercado laboral \cite{cepal202X}. No obstante, esta expansión ha estado acompañada de una notable proliferación en la oferta académica nacional \cite{sies2024informe}, evidenciada por una multiplicidad de nomenclaturas institucionales, tales como Ingeniería Civil Informática e Ingeniería en Informática, entre otras. Esta diversidad nominal sugiere una heterogeneidad intrínseca en los contenidos formativos y en los perfiles de egreso \cite{mineduc2021propuesta}, lo que genera un escenario de incertidumbre respecto a la calidad, el enfoque y la pertinencia de la educación impartida. En la actualidad, se evidencia la carencia de una caracterización sistemática y actualizada de estos programas de estudio \cite{mineduc2021propuesta}, lo que dificulta la comprensión objetiva de las diferencias sustanciales entre las distintas opciones académicas, particularmente en métricas cuantitativas de similitud entre mallas curriculares y de alineación entre perfiles de egreso y contenidos formativos.
        %--------------------------------------------------
        % C.1.2) Contexto
        %--------------------------------------------------
    

\subsection{Contexto} \label{Contexto}
En las últimas décadas, el acelerado desarrollo tecnológico y la creciente digitalización de los procesos industriales han transformado profundamente el sector productivo en Chile, incrementando la demanda por competencias técnicas avanzadas y especializadas \cite{ITTalent2021}. Como respuesta, el ecosistema educativo nacional ha experimentado una diversificación significativa, caracterizada por una proliferación de programas con denominaciones heterogéneas y una oferta que se ha adaptado mediante cambios en las modalidades de estudio —como el aumento de programas vespertinos y no presenciales— y ajustes en la duración de las carreras para mejorar la eficiencia terminal \cite{LANDONI2015}.

Es en este escenario que se manifiesta la problemática fundamental: la inexistencia de una caracterización sistemática y actualizada a nivel nacional que permita comparar de forma objetiva la diversidad de los programas vigentes. Si bien el sistema ha transitado desde un modelo rígido hacia una mayor flexibilidad curricular \cite{LANDONI2015}, la falta de descriptores sustantivos sobre los resultados de aprendizaje genera una "opacidad" en las credenciales académicas \cite{LANDONI2015}. Aunque existen marcos de referencia internacionales como los lineamientos de la IEEE y la ACM, utilizados globalmente para la definición curricular \cite{CC2020TaskForce2021}, el esfuerzo nacional más cercano ha sido el portal Mi Futuro \cite{mifuturo_buscador}, el cual, si bien provee datos sobre tendencias laborales y empleabilidad, no constituye una caracterización integral que permita evaluar comparativamente la heterogeneidad curricular entre instituciones y programas.
        %--------------------------------------------------
        % C.1.3) Formulación del problema
        %--------------------------------------------------
\subsection{Formulación del problema} \label{formulación}

Dado el contexto anterior, surge la siguiente interrogante principal: ¿Cómo se alinean los perfiles formativos de los programas de Ingeniería Informática en Chile con los estándares internacionales IEE/ACM y las competencias demandadas por el mercado laboral en el sector TI?  
        %--------------------------------------------------
        % C.1.4) Objetivo preliminar
        %--------------------------------------------------

\subsection{Objetivo preliminar} \label{objetivo}
Diseñar una metodología de análisis comparativo para clasificar y comparar los perfiles curriculares de Ingeniería Informática en Chile, mediante Procesamiento de Lenguaje Natural (NLP) y técnicas de clustering, con el fin de medir la brecha de competencias respecto de los estándares internacionales IEEE/ACM y los requerimientos del mercado laboral.
        %--------------------------------------------------
        % C.1.5) Propuesta de solución
        %--------------------------------------------------
\subsection{Propuesta} \label{Propuesta}

Para abordar la falta de caracterización sistemática de estos programas, se propone diseñar un modelo de análisis estructural de datos y Procesamiento de Lenguaje Natural (NLP). Este enfoque procesará tres dimensiones multidimensionales clave de la oferta académica: los perfiles de egreso, las estructuras de las mallas curriculares y los perfiles investigativos declarados por las instituciones.

La metodología contempla una fase de extracción de características y vectorización semántica, sobre la cual se aplicarán algoritmos de agrupamiento no supervisado. El agrupamiento operará sobre variables específicas, tales como la distribución de áreas de conocimiento, la frecuencia de descriptores de competencias técnicas y la densidad de enfoques de investigación.

Esta segmentación paramétrica permitirá calcular la distancia geométrica o semántica para contrastar la oferta académica frente a las demandas reales del mercado laboral y los estándares internacionales (IEEE/ACM). Finalmente, los clústeres identificados y las métricas de brecha se integrarán en una herramienta de visualización interactiva, habilitando la exploración dinámica y la interpretación integral de los hallazgos.
        %--------------------------------------------------
        % C.1.6) Justificación
        %--------------------------------------------------
        
\subsection{Justificación} \label{Justificación}

La justificación de esta investigación radica en la necesidad imperativa de superar los diagnósticos curriculares tradicionales, los cuales suelen ser manuales y de lenta adaptación. Actualmente, la carencia de una caracterización estructurada dificulta la alineación objetiva entre la formación académica y las competencias técnicas específicas exigidas por el mercado laboral TI\cite{ITTalent2021}.

El aporte fundamental de este trabajo trasciende la mera aplicación de técnicas de Procesamiento de Lenguaje Natural y algoritmos de agrupamiento; su verdadero valor reside en la generación de conocimiento nuevo al establecer una línea base empírica y cuantitativa sobre el estado de la Ingeniería Informática en Chile..

El impacto esperado de este modelo analítico es dotar de transparencia al sistema educativo, entregando una herramienta fundamentada en evidencia real. Esto facilitará la toma de decisiones estratégicas continuas para reducir las brechas de competencias y asegurar la pertinencia de los perfiles formativos frente a los estándares de calidad internacionales.

        %--------------------------------------------------
        % C.1.7) Viabilidad
        %--------------------------------------------------
\subsection{Viabilidad} \label{Viabilidad}

La viabilidad del proyecto se sustenta en el acceso garantizado a fuentes de información pública provenientes de portales institucionales de las universidades chilenas y repositorios gubernamentales como 'Mi Futuro'\cite{mifuturo_buscador}, los cuales consolidan las universidades que ofrecen el programa de Ingeniería Informática.  A nivel técnico, se identifica que la disponibilidad de los datos se encuentra mayoritariamente en formatos no estructurados, tales como documentos PDF y estructuras HTML. Por consiguiente, el estudio contempla un pipeline de preprocesamiento crítico que representa aproximadamente el 40 por ciento del esfuerzo operativo, enfocado en la extracción de texto (scraping y parsing), la limpieza de ruido semántico y la estandarización de nomenclaturas heterogéneas.  Operativamente, la investigación es factible mediante el uso de herramientas de código abierto para el procesamiento de grandes volúmenes de texto y la ejecución de modelos de representación vectorial. Se cuenta con los recursos computacionales necesarios para el entrenamiento y despliegue de los algoritmos de agrupamiento y la plataforma de visualización.


%--------------------------------------------------
% C.X) Bibliography
%--------------------------------------------------    
\bibliographystyle{IEEEtran}
\bibliography{REF/biblio}

\end{document}